\section{Modelação de Intensidade e Processos de Contagem}

Para modelar eventos recorrentes, abandonamos a função de sobrevivência simples $S(t)$ em favor da \textbf{Função de Intensidade} $\lambda(t | \mathcal{F}_{t^-})$, onde $\mathcal{F}_{t^-}$ representa a história da empresa até ao instante imediatamente anterior a $t$.

\subsection{Modelo de Prentice-Williams-Peterson (PWP-GT)}
Ao contrário do modelo de Andersen-Gill, o PWP permite que o risco mude com o número de episódios ($k$), sendo o modelo de eleição para testar a nossa hipótese H3:

\begin{equation}
    \lambda_k(t | Z) = \lambda_{0k}(t) \exp(\beta_k' Z)
\end{equation}

Onde:
\begin{itemize}
    \item $\lambda_{0k}(t)$ é o hazard base para o $k$-ésimo episódio de stress.
    \item $\beta_k$ representa os coeficientes específicos para o episódio $k$, permitindo verificar se a importância da \textit{Governance} ou \textit{Macro} se altera após o primeiro colapso.
\end{itemize}

\section{Riscos Concorrentes (Competing Risks)}

No contexto de PME, o evento terminal (Falência) atua como um risco concorrente para a recuperação. No modelo \textit{Multi-State}, a transição do Estado 1 (Distress) tem dois caminhos mutuamente exclusivos:
\begin{enumerate}
    \item \textbf{Transição $1 \to 0$:} Recuperação para a saúde financeira.
    \item \textbf{Transição $1 \to 2$:} Progressão para a falência total (Estado Absorvente).
\end{enumerate}

A função estatística correta para analisar este fenómeno é a \textbf{Função de Incidência Cumulativa (CIF)}:
\begin{equation}
    CIF_k(t) = P(T \le t, \text{Cause} = k)
\end{equation}
Esta função é superior ao estimador de Kaplan-Meier nestes cenários, pois não assume independência entre as causas de falha.

\section{Arquitetura Dinâmica do DeepHit}
\todo[inline,color=red!40]{PENDENTE: Descarregar Lee (2019) Dynamic-DeepHit para /articles}

O \textbf{DynamicDeepHit} estima a probabilidade discreta de falha para cada causa $k$ e cada intervalo de tempo $t$, utilizando a trajetória histórica $X_{1:T}$:
\begin{equation}
    y = f_{\theta}(X_{1:T})
\end{equation}
A rede neuronal é composta por uma \textbf{RNN/LSTM partilhada} que captura a dinâmica temporal, seguida de \textbf{MLPs causa-específicas} que estimam o risco para Distress vs. Falência.
